% Forza in funzione del tempo durante l'arresto di un passeggero: urto contro una
% superficie rigida (breve e intenso) e con l'airbag (lungo e meno intenso).
% Le aree, cioè gli impulsi, sono uguali.
\begin{tikzpicture}[font=\small]
  \begin{axis}[
    width=9cm, height=5.2cm,
    axis lines=left, xmin=0, xmax=1.05, ymin=0, ymax=10.5,
    xlabel={$t$}, ylabel={$F$}, xtick=\empty, ytick=\empty,
    every axis x label/.style={at={(current axis.right of origin)}, anchor=west},
    every axis y label/.style={at={(current axis.above origin)}, anchor=south},
    clip=false]
    \addplot[attenzione, very thick, fill=attenzione!20, fill opacity=0.6, domain=0:0.4, samples=100]
      {9*exp(-((x-0.2)/0.05)^2)} \closedcycle;
    \addplot[formula, very thick, fill=formula!20, fill opacity=0.6, domain=0:1.0, samples=100]
      {1.5*exp(-((x-0.5)/0.3)^2)} \closedcycle;
    \node[attenzione, anchor=west, align=left] at (axis cs:0.27,8) {superficie rigida:\\$\Delta t$ breve, forza grande};
    \node[formula, anchor=south, align=center] at (axis cs:0.72,1.7) {airbag: $\Delta t$ lungo,\\forza piccola};
  \end{axis}
\end{tikzpicture}
